% ---------------------------------------------------------------
% Section 3: Preliminaries
% ---------------------------------------------------------------

\section{Preliminaries}
\label{sec:preliminaries}

We establish the mathematical foundations for our diagnostic analysis, covering contrastive vision-language learning, the BioCLIP model family, Linnaean taxonomy, and the embedding geometry metrics used throughout.

\subsection{Contrastive Vision-Language Learning}

\noindent\textbf{CLIP-style Training.}
Given image-text pairs $\{(x_i, t_i)\}_{i=1}^{N}$, CLIP learns $\ell_2$-normalized encoders
$f_v : \mathcal{X} \to \mathbb{R}^{d}$ and $f_t : \mathcal{T} \to \mathbb{R}^{d}$
by minimizing a symmetric InfoNCE loss over mini-batches of size $B$:

\begin{equation}
\mathcal{L}_{\mathrm{CLIP}}
= -\frac{1}{2B}\sum_{i=1}^{B}
\left[
\log\frac{e^{\mathrm{sim}(x_i,t_i)/\tau}}{\sum_{j} e^{\mathrm{sim}(x_i,t_j)/\tau}}
+
\log\frac{e^{\mathrm{sim}(x_i,t_i)/\tau}}{\sum_{j} e^{\mathrm{sim}(x_j,t_i)/\tau}}
\right],
\label{eq:clip_loss}
\end{equation}

where $\mathrm{sim}(x,t)=f_v(x)^{\top}f_t(t)$ and $\tau>0$ is a learned temperature.
At convergence, semantically matched pairs cluster together in the shared space.

\noindent\textbf{Zero-shot Classification.}
At inference, text prompts $\{t_c\}_{c=1}^{K}$ are encoded into class prototypes
$\{w_c = f_t(t_c)\}$, and a test image $x$ is assigned to
$\hat{c} = \arg\max_{c}\,\mathrm{sim}(x, t_c)$.
The choice of prompt $t_c$ directly shapes the geometry of class prototypes and
nearest-neighbor decision boundaries, making prompt design the central object of
this work.

\noindent\textbf{Alignment and Uniformity.}
Wang and Isola~\cite{wang2020understanding} characterize embedding quality via
\emph{alignment} (cohesion of positive pairs) and \emph{uniformity} (spread across
the hypersphere). We report these as auxiliary global-geometry summaries alongside
the cluster-level metrics in Section~\ref{sec:geometry_metrics}.

\subsection{BioCLIP and BioCLIP2}

\noindent\textbf{BioCLIP}~\cite{stevens2024bioclip} introduced hierarchical taxonomic
prompts to biological vision-language learning, prepending the full Linnaean lineage
(e.g., \texttt{Animalia Chordata Aves \ldots Passer domesticus}) to yield 17--20 pp
gains over OpenCLIP. Qualitative analysis suggests the embedding space conforms to
the tree of life, but no formal geometric characterization is given.

\noindent\textbf{BioCLIP2}~\cite{gu2025bioclip2} scales this paradigm to
TreeOfLife-200M ($\approx$214M images, ${>}$500K taxa) with a ViT-L/14 image encoder
and an autoregressive text encoder, applying a separate contrastive loss at each rank.
It achieves $+18.1$ pp over BioCLIP and exhibits \emph{emergent properties}:
inter-species embeddings align with ecological traits (beak size, habitat) while
intra-species variation is preserved in orthogonal subspaces.
These properties are documented qualitatively; the present work provides their first
controlled quantitative characterization.

\noindent\textbf{Species-level Saturation.}
Training on TreeOfLife-200M yields near-ceiling species discrimination \emph{before}
any hierarchical text is applied.
On CUB-200-2011 (200 bird species, 11,788 images), the flat-prompt baseline already
achieves a species-level silhouette of $0.775$ and kNN purity of $99.9\%$.
In this \emph{saturation regime}, hierarchical text cannot improve already near-perfect
species boundaries and may instead dilute them by pulling embeddings toward shared
higher-rank anchors — a key interpretive frame for the results in Section~\ref{sec:experiments}.

\subsection{Linnaean Taxonomic Hierarchy}

\noindent\textbf{Rank Structure.}
We adopt the standard seven-rank Linnaean system ordered from coarsest to finest:
\begin{equation}
\text{kingdom} \succ \text{phylum} \succ \text{class} \succ \text{order}
\succ \text{family} \succ \text{genus} \succ \text{species}.
\label{eq:linnaean_ranks}
\end{equation}
Each species $c$ has a unique lineage $\ell(c) = (c^{(1)}, \dots, c^{(7)})$;
higher-rank taxa are shared across species within the same clade.

\noindent\textbf{Taxonomic Proximity.}
The \emph{lowest common ancestor} (LCA) of two species is the deepest rank at which
their lineages agree; its depth in $\{1,\dots,7\}$ is our ground-truth measure of
taxonomic closeness, underlying the latent-taxonomy probe in
Section~\ref{sec:rq2}.

\noindent\textbf{Rank Homogeneity Confound.}
When all species in a dataset share the same taxon at rank $r^{*}$, the corresponding
prompt token is constant across every prompt and carries zero discriminative
information, acting as uninformative anchor noise.
CUB-200-2011 exemplifies this: all 200 species are Aves, so
\texttt{Animalia}, \texttt{Chordata}, and \texttt{Aves} are identical across every
prompt, effectively reducing the seven-rank hierarchy to four ranks
(order $\to$ species).
We address this confound by replicating experiments on the Rare Species dataset
(400 species spanning multiple phyla), where upper-rank tokens are taxonomically
diverse (Figure~\ref{fig:rank_homogeneity}).

\begin{figure}[t]
  \centering
  % \includegraphics[width=0.95\linewidth]{figures/rank_homogeneity.pdf}
  \caption{Rank homogeneity confound.
           \textbf{Left (CUB-200-2011):} all 200 species share kingdom, phylum, and
           class (Aves), so the top three rank tokens (grey) are constant and
           contribute no inter-species signal.
           \textbf{Right (Rare Species):} 400 species span multiple phyla, enabling
           evaluation across all seven ranks.
           Coloured nodes indicate rank tokens that vary across species;
           grey nodes indicate tokens shared by all species in the dataset.}
  \label{fig:rank_homogeneity}
\end{figure}

\subsection{Embedding Geometry Metrics}
\label{sec:geometry_metrics}

Let $Z_c = \{z_i : y_i^{(\mathrm{species})} = c\}$ be the set of $\ell_2$-normalized
image embeddings for species $c$, with unnormalized centroid $\mu_c$ and normalized
centroid $\hat{\mu}_c = \mu_c / \|\mu_c\|$.

\noindent\textbf{Intra-class Variance} ($\downarrow$).
Mean pairwise cosine distance within each species, macro-averaged over all species:
\begin{equation}
\overline{\sigma^2}
= \frac{1}{|\mathcal{C}|}
  \sum_{c}\;
  \frac{1}{\binom{|Z_c|}{2}}
  \sum_{\substack{z,z'\in Z_c \\ z\neq z'}}
  (1 - z^{\top}z').
\label{eq:intra_var}
\end{equation}

\noindent\textbf{Inter-class Margin} ($\uparrow$).
Nearest-centroid distance normalized by within-class spread, macro-averaged:
\begin{equation}
\overline{\delta}
= \frac{1}{|\mathcal{C}|}
  \sum_{c}
  \frac{\min_{c'\neq c}(1 - \hat{\mu}_c^{\top}\hat{\mu}_{c'})}
       {\sigma_c}.
\label{eq:inter_margin}
\end{equation}

\noindent\textbf{Silhouette Score at Rank $r$} ($\uparrow$).
Our \emph{primary} metric. Let $a_i^{(r)}$ and $b_i^{(r)}$ be the mean intra- and
nearest-inter-cluster cosine distances for $z_i$ at rank $r$:
\begin{equation}
S^{(r)}
= \frac{1}{N}
  \sum_{i=1}^{N}
  \frac{b_i^{(r)} - a_i^{(r)}}
       {\max(a_i^{(r)},\, b_i^{(r)})}.
\label{eq:silhouette}
\end{equation}
Computed at each of the seven ranks, $S^{(r)}$ quantifies geometric organization at
every level of the taxonomy.

\noindent\textbf{kNN Purity at Rank $r$} ($\uparrow$).
Fraction of $k{=}10$ nearest neighbors sharing the same rank-$r$ label:
\begin{equation}
\mathrm{Purity}^{(r)}_k
= \frac{1}{N}
  \sum_{i=1}^{N}
  \frac{|\{j \in \mathrm{NN}_k(z_i) : y_j^{(r)} = y_i^{(r)}\}|}{k}.
\label{eq:knn_purity}
\end{equation}

\noindent\textbf{RankMe} ($\uparrow$).
Effective rank of the $N{\times}d$ embedding matrix~\cite{garrido2023rankme}:
\begin{equation}
\mathrm{RankMe}(\mathbf{Z})
= \exp\!\Bigl(-\textstyle\sum_{k} p_k \log p_k\Bigr),
\quad
p_k = \frac{\sigma_k(\mathbf{Z})}{\|\sigma(\mathbf{Z})\|_1},
\label{eq:rankme}
\end{equation}
where $\sigma(\mathbf{Z})$ is the vector of singular values of $\mathbf{Z}$.

All metrics are computed on $\ell_2$-normalized embeddings. We report
\emph{image-only} variants ($\{z_i\}$ alone) alongside \emph{image-text} variants
(with text prototypes $\{w_c\}$ as anchors) to isolate modality-gap
artefacts~\cite{liang2022mindgap} from genuine geometric effects.
